\begin{enumerate}

\item \textbf{การหาปริพันธ์ของฟังก์ชันเศษส่วนพหุนาม}
จงหาปริพันธ์ต่อไปนี้
\begin{tasks}(2)
\task $\int \frac{x^2+1}{x} \, dx$
\vspace{150pt}
\task $\int \frac{x^2+1}{x+1} \, dx$
\end{tasks}

\item \textbf{การหาแปลงเป็นเศษส่วนย่อย}
จงเขียนฟังก์ชันต่อไปนี้ ในรูปเศษส่วนย่อย
\begin{tasks}(2)
\task $\frac{1}{(x-1)(2x-1)}$
\vspace{150pt}
\task $\frac{x+1}{(x+3)(x+2)}$
\end{tasks}


\item \textbf{การหาปริพันธ์โดยใช้เศษส่วนย่อย}
กำหนดสูตรปริพันธ์ดังนี้
\[ \int \frac{1}{rx+s} \, dx = \rule{100pt}{1pt}\]
จงหาปริพันธ์ต่อไปนี้ โดยการแยกเป็นเศษส่วนย่อย
\begin{tasks}(2)
\task $\int \frac{1}{(x-1)(2x-1)} \, dx$
\vspace{100pt}
\task $\int \frac{x+1}{(x+3)(x+2)} \, dx$
\end{tasks}

\newpage

\item \textbf{การหาปริพันธ์ด้วยการแทนค่ามุมครึ่ง}

กำหนดให้ $u = \tan \frac{x}{2}$ แล้วจะได้
\begin{align*}
\sin x &= \rule{100pt}{1pt} \\
\cos x &= \rule{100pt}{1pt} \\
dx &= \rule{100pt}{1pt}
\end{align*}

จงหาปริพันธ์ต่อไปนี้ โดยการแทนค่ามุมครึ่ง
\begin{tasks}[after-item-skip = 200pt, after-skip=200pt](2)
\task $\int \frac{1}{1 + \cos x} \, dx$
\task $\int \frac{1}{\sin x} \, dx$
\end{tasks}


\item \textbf{การหาปริพันธ์ด้วยการแทนค่าด้วยตัวแปรอื่น}

จงหาปริพันธ์ต่อไปนี้ โดยการแทนค่าด้วยตัวแปรอื่น
\begin{tasks}[after-item-skip = 100pt, after-skip=100pt](2)
\task $\int \cos (\sqrt{x}) \, dx$
\task $\int \frac{1}{2+\sqrt{x-4}} \, dx$
\end{tasks}

\end{enumerate}